\subsection{Relaxation Cascade, Dynamical $\Lambda_{\mathrm{eff}}$,
  and the Cosmological Origin of Fundamental Constants}
  \label{sec:relax-cascade}

  The preceding sections establish $\chi(t)$ as the global relaxation field
  and $H(t) = \dot{\chi}/\chi$ as the emergent Hubble rate (Section~7.4,
  Eq.~(67)).
  This section develops four structural consequences that follow directly
  from these definitions together with the constitutive relations
  Eqs.~(202) and~(204):
  \begin{enumerate}
    \item the effective cosmological constant $\Lambda_{\mathrm{eff}}$ is
    not a fundamental constant but a dynamical residual of ongoing
    $\chi$-relaxation, with equation of state $w \neq -1$;
    \item relaxation proceeds as a discrete cascade of projectable
    configurations, providing a natural explanation for the
    $10^{16}$--$10^{19}$ hierarchy between the Planck energy and
    observed particle energies;
    \item Planck's constant $\hbar$ is an epoch-dependent invariant of
    the current relaxation regime, fixed by $K_0$ and $\chi_c$ through
    Eq.~(202);
    \item the Big Bang corresponds precisely to the first globally
    projectable configuration of $\chi$, whose existence is
    characterised by a finite saturation condition derived from the
    Born--Infeld threshold.
  \end{enumerate}
  Each result is an internal consequence of the Cosmochrony framework;
  no new free parameter is introduced.

  \subsubsection*{Dynamical $\Lambda_{\mathrm{eff}}$ from Residual
  Relaxation}
    \label{sec:relax-cascade-lambda}

    \paragraph{Definition of residual relaxation energy.}
      Energy is defined in this framework as the operational cost of modifying
      a projected configuration while preserving admissibility
      (Appendix~F.7).
      At the global cosmological level, the outstanding relaxation capacity
      of $\chi$ at cosmic time $t$ is the quantity that remains to be
      released before $\chi$ reaches an asymptotically equilibrated state.
      We define it as a dimensionless fractional quantity $\Omega_\chi(t)$
      already introduced in Eq.~(168):
    \begin{equation}
      \label{eq:rc-omega}
      \Omega_\chi(t) \;\equiv\; \langle \beta^2 \rangle,
      \qquad
      \beta \;\equiv\; \frac{|\nabla\chi|}{c},
      \end{equation}
      which measures the fraction of relaxation capacity stored in active
      spatial gradients.
      The effective cosmological density associated with residual large-scale
      relaxation (Appendix~F.6) is then
    \begin{equation}
      \label{eq:rc-rho-lambda}
      \rho_{\Lambda}(t)
      \;\equiv\;
      \frac{3\,\Omega_\chi(t)\,H^2(t)}{8\pi G},
      \end{equation}
      which is dimensionally consistent with the critical density
      $\rho_c = 3H^2/(8\pi G)$ (natural units with $c=1$ at the effective
      spacetime level; factors of $c$ are restored in numerical estimates
      below).
      In the homogeneous background limit where $\Omega_\chi$ is nearly
      constant, Eq.~\eqref{eq:rc-rho-lambda} reproduces the Friedmann
      closure condition up to an $O(\Omega_\chi)$ correction, consistently
      with Appendix~C.3, Eq.~(169).

    \paragraph{Derivation of $\rho_\Lambda \propto H^2$.}
      From the linear background closure $\chi(t) = \chi_0 + ct$ (Eq.~(66))
      and the definition $H(t) = \dot\chi/\chi$ (Eq.~(67)):
    \begin{equation}
      \label{eq:rc-H-explicit}
      H(t) \;=\; \frac{c}{\chi_0 + ct}
      \;\xrightarrow{t \gg \chi_0/c}\; \frac{1}{t},
      \end{equation}
      so $H \propto t^{-1}$ holds in the late-time relaxation-dominated
      regime.
      Since $\rho_\Lambda \propto \Omega_\chi H^2$ and $\Omega_\chi$ varies
      slowly (Appendix~C.3), one obtains directly
    \begin{equation}
      \label{eq:rc-rho-H2}
      \rho_\Lambda(t)
      \;\propto\; H^2(t)
      \;\propto\; \frac{1}{t^2}.
      \end{equation}
      Numerically, with $H_0 \approx 2.2 \times 10^{-18}~\mathrm{s}^{-1}$
      and $G = 6.67 \times 10^{-11}~\mathrm{m^3\,kg^{-1}\,s^{-2}}$,
    \begin{equation}
      \label{eq:rc-rho-num}
      \rho_\Lambda(t_0)
      \;\sim\;
      \Omega_\chi(t_0)\,\frac{3 H_0^2}{8\pi G}
      \;\approx\;
      \Omega_\chi(t_0)\times 9 \times 10^{-27}~\mathrm{kg\,m^{-3}},
      \end{equation}
      in agreement with the observed dark energy density for
      $\Omega_\chi(t_0) \sim 0.7$, without any fine-tuning.

    \paragraph{Equation of state $w \neq -1$.}
      With $\Omega_\chi$ treated as a slowly varying function of the scale
      factor $a$, differentiate $\rho_\Lambda = 3\Omega_\chi H^2/(8\pi G)$
      along the background:
    \begin{equation}
      \label{eq:rc-rhodot}
      \frac{\dot\rho_\Lambda}{\rho_\Lambda}
      \;=\; \frac{\mathrm{d}\ln\Omega_\chi}{\mathrm{d}\ln a}\,H
      \;+\; 2\frac{\dot H}{H}.
      \end{equation}
      Substituting into the continuity equation
      $\dot\rho_\Lambda + 3H(1+w_\chi)\rho_\Lambda = 0$ yields
    \begin{equation}
      \label{eq:rc-w}
      w_\chi
      \;=\; -1 + \frac{2(1+q)}{3}
      \;-\; \frac{1}{3}\frac{\mathrm{d}\ln\Omega_\chi}{\mathrm{d}\ln a},
      \end{equation}
      where $q = -\ddot{a}/(aH^2)$ is the deceleration parameter of the
      \emph{same} self-consistent background, and all three terms are
      dynamical outputs of the framework.

      Two remarks are essential.
      First, the correction term $-\tfrac{1}{3}\,\mathrm{d}\ln\Omega_\chi/\mathrm{d}\ln a$
      is not a free parameter: $\Omega_\chi(a) = \langle\beta^2\rangle$
      (Eq.~\eqref{eq:rc-omega}) is determined by the relaxation dynamics.
      Its contribution can have either sign.
      If $\Omega_\chi$ decreases as $\Omega_\chi \propto a^{-\delta}$ with
      $\delta > 0$, then
      $-\tfrac{1}{3}\,\mathrm{d}\ln\Omega_\chi/\mathrm{d}\ln a = +\delta/3$,
      which increases $w_\chi$ (i.e.\ makes it less negative).
      Conversely, any epoch where $\Omega_\chi$ grows with $a$ contributes
      negatively and can keep $w_\chi$ closer to $-1$ even when $q$ departs
      from $-1$.

      Second, using the $\Lambda$CDM-inferred value $q \approx -0.55$
      directly in Eq.~\eqref{eq:rc-w} would be circular: that value of $q$
      is reconstructed assuming a dark energy component with $w = -1$,
      whereas the present model changes the background equations.
      The proper observational confrontation is therefore not through a
      single effective $w$, but through the predicted $H(z)$ profile
      (Appendix~C.3, Eq.~(166)), which can be compared directly to
      BAO and supernova data without assuming a fixed equation of state.
      The framework's prediction of a mild $H(z)$ enhancement at
      $0.1 \lesssim z \lesssim 10$ relative to $\Lambda$CDM
      (Section~11.1) is the appropriate discriminating observable.

    \paragraph{Resolution of the coincidence problem.}
      In the standard $\Lambda$CDM framework, the coincidence between
      $\rho_\Lambda$ and the matter density $\rho_m$ today is accidental.
      Here, both $\rho_\Lambda \propto H^2$ and $\rho_m \propto H^2$ (in the
      matter-dominated limit $\rho_m = 3H^2/(8\pi G)$), so their ratio
      $\rho_\Lambda / \rho_m \sim \Omega_\chi$ remains of order unity
      throughout the relaxation-dominated epoch.
      The coincidence is therefore a structural consequence of the single
      relaxation dynamics driving both matter inertia and residual vacuum
      energy.

  \subsubsection*{Discrete Relaxation Cascade and the Energy Hierarchy}
    \label{sec:relax-cascade-hier}

    \paragraph{Discreteness of projection.}
      The projectability threshold $\Theta_p$ (Appendix~D.5) defines the
      minimal spectral gap above which a configuration of $\chi$ admits a
      stable geometric description.
      Each time the relaxation network crosses $\Theta_p$ from above, a new
      projectable regime opens: this is what we call a \emph{relaxation step}
      $n \to n+1$.
      The minimum action exchanged per step is bounded below by the
      projection granularity
    \begin{equation}
      \label{eq:rc-delta-S}
      \Delta S_n \;\geq\; \hbar_\chi
      \;\equiv\;
      \frac{c^3}{K_0\,\chi_c},
      \end{equation}
      which is Eq.~(202).
      Relaxation is therefore intrinsically discrete: the continuum effective
      description is an approximation valid between successive steps.

    \paragraph{Cascade law.}
      Let $E_n$ denote the effective energy available to the projected
      description at step $n$.
      Since each step irreversibly increases the number of distinguishable
      projected configurations $\Omega_n$, the energy per configuration
      decreases.
      Under the assumption that each step redistributes the remaining energy
      uniformly over the newly accessible configurations, and using the
      self-similar spectral structure established in Appendix~D.2
      (the ratio $G_\mathrm{eff} \propto K_0/\chi_c^2$ is preserved under
      renormalisation), one obtains
    \begin{equation}
      \label{eq:rc-cascade}
      E_n \;\sim\; \frac{E_P}{n},
      \end{equation}
      where $E_P = \sqrt{\hbar c^5/G} \approx 2 \times 10^9~\mathrm{J}$
      is the Planck energy, corresponding to the saturated initial state
      (Section~7.1).

    \paragraph{Total number of steps and the hierarchy.}
      The total number of Planck-time steps elapsed since the initial
      projection is
    \begin{equation}
      \label{eq:rc-N}
      N \;\equiv\; \frac{t_U}{t_P}
      \;\approx\; \frac{4.3 \times 10^{17}~\mathrm{s}}{5.4 \times 10^{-44}~\mathrm{s}}
      \;\approx\; 10^{61}.
      \end{equation}
      Applying Eq.~\eqref{eq:rc-cascade} with $n = N$:
    \begin{equation}
      \label{eq:rc-Eparticle}
      E_\mathrm{particle}
      \;\sim\; \frac{E_P}{N}
      \;\sim\; \frac{2 \times 10^9}{10^{61}}~\mathrm{J}
      \;\sim\; 10^{-52}~\mathrm{J}
      \;\sim\; 10^{-33}~\mathrm{eV},
      \end{equation}
      which corresponds to the energy scale associated with the cosmological
      constant (the dark energy quantum $\hbar H_0$), consistently with
      Eq.~\eqref{eq:rc-rho-num}.
      For intermediate particle species stabilised at earlier steps
      $n \ll N$, Eq.~\eqref{eq:rc-cascade} yields
    \begin{equation}
      \label{eq:rc-hierarchy}
      \frac{E_P}{E_\mathrm{particle}}
      \;\sim\; n,
      \end{equation}
      so the observed hierarchy $E_P/m_e c^2 \approx 10^{22}$ requires
      that the electron stabilised at a relaxation step
      $n_e \sim 10^{22}$, far earlier than the present epoch.
      This provides a structural interpretation of the hierarchy problem:
      particle masses reflect the depth of their stabilisation step in the
      relaxation cascade, not an unexplained fine-tuning of a Higgs
      potential.

    \paragraph{Primordial energy and cascade consistency.}
      Inverting Eq.~\eqref{eq:rc-cascade} at step $n = 1$, the energy
      available at the first projectable configuration is $E_0 \sim E_P$,
      in agreement with the maximal-constraint interpretation of the Big Bang
      (Section~7.1).
      A crude upper bound on the total observable energy today follows from
      counting all Planck-volume cells in the observable Universe,
      $N_\mathrm{vol} \sim (R_U/\ell_P)^3 \sim 10^{183}$, and assigning
      each cell the energy $E_P/N \sim 10^{-52}~\mathrm{J}$ of the current
      cascade step.
      This bound far exceeds the observed $E_\mathrm{today} \sim 10^{69}~\mathrm{J}$
      because it ignores the spectral distribution of stabilisation depths:
      most cells have relaxed to states with $n \gg N$ and contribute
      negligibly.
      A quantitative account requires the full spectral distribution of
      stabilisation depths across the relaxation network, which is left for
      future work.
      The cascade law \eqref{eq:rc-cascade} is therefore presented as a
      structural framework for the hierarchy, not as a closed energy budget.

  \subsubsection*{Planck's Constant as a Relaxation-Epoch Invariant}
    \label{sec:relax-cascade-hbar}

    \paragraph[Starting point: Eq. hbar-chi-def]{Starting point: Eq.~\eqref{eq:hbar_chi_def}.}
      The effective action quantum is defined structurally as
    \begin{equation}
      \label{eq:rc-hbar-def}
      \hbar_\chi
      \;\equiv\;
      \frac{c^3}{K_0\,\chi_c},
      \end{equation}
      where $K_0$ is the maximal relaxation conductivity and $\chi_c$ is the
      critical relaxation threshold (Appendix~F.1).
      The identification $\hbar_\chi \simeq \hbar$ is operationally enforced
      in projectable regimes (Appendix~D.6, Eq.~(185)).

    \paragraph{Cosmological normalisation.}
      Appendix~D.2 establishes that the dynamics are spectrally self-similar:
      $G$ remains constant under renormalisation provided the ratio
      $K_0/\chi_c^2$ is preserved.
      The cosmological-scale identification $\chi_c \sim c/H_0 = \chi(t_0)$
      (Appendix~C.4) then gives, using Eq.~(204),
    \begin{equation}
      \label{eq:rc-K0-cosmo}
      K_0 \;\sim\; \frac{c^4}{16\pi G\,\chi_c^2}
      \;\sim\; \frac{H_0^2}{16\pi G}
      \;\sim\; 10^{-52}~\mathrm{m}^{-2},
      \end{equation}
      which is the cosmological-scale entry of Table~4.
      Substituting into Eq.~\eqref{eq:rc-hbar-def}:
    \begin{equation}
      \label{eq:rc-hbar-cosmo}
      \hbar_\chi
      \;\sim\; \frac{c^3}{K_0\,\chi_c}
      \;\sim\; \frac{c^3\,\cdot\, 16\pi G}{H_0^2/H_0}
      \;\sim\; \frac{16\pi G c^3}{H_0}
      \;\approx\; 10^{70}~\mathrm{J\,s},
      \end{equation}
      a cosmological action scale, not the microscopic $\hbar$.
      The microscopic Planck-scale identification $\chi_c \sim \ell_P$
      instead yields $K_0 \sim 10^{93}~\mathrm{m}^{-2}$ and recovers
      $\hbar_\chi \simeq \hbar = 1.05 \times 10^{-34}~\mathrm{J\,s}$.
      Both branches are internally consistent, as stated in Appendix~D.2.

    \paragraph{Interpretation.}
      The two branches reflect the spectral self-similarity of the cascade:
      at each scale, the ratio $K_0\,\chi_c^2$ is fixed by $G$, while
      $\hbar_\chi$ tracks the granularity of that scale's projection.
      The universal microscopic $\hbar$ is therefore the action granularity
      of the \emph{current particle-scale} projectable regime.
      It is not a fundamental constant of $\chi$; it is the action quantum
      that emerges when the projection resolves configurations at the
      length scale $\ell_P$, as encoded by the Planck-branch normalisation.
      The cosmological branch ($\hbar_\chi \sim 10^{70}~\mathrm{J\,s}$)
      corresponds to the action scale of the relaxation field itself on
      Hubble scales and is not directly observable in particle dynamics;
      it characterises the global relaxation flux, not the granularity of
      quantum transitions in the effective projected description.
      This implies a further structural consequence: in primordial
      high-constraint regimes approaching the initial projection
      (Section~7.1), $\chi_c$ may differ from its present-epoch value,
      causing both $\hbar$ and the fine-structure constant $\alpha$ to scale
      accordingly (Eq.~(9)), while leaving the product
      $K_0\,\chi_c^2 \propto G^{-1}$ invariant.

  \subsubsection*{The Big Bang as the First Globally Projectable
  Configuration}
    \label{sec:relax-cascade-bigbang}

    \paragraph{Projectability condition.}
      A configuration of $\chi$ is projectable if and only if its local
      Born--Infeld saturation parameter (Appendix~D.5, Eq.~(179)) satisfies
    \begin{equation}
      \label{eq:rc-proj-cond}
      S_i(\chi) \;\equiv\;
      \frac{1}{c^2}\sum_{j \in \mathcal{N}(i)} K_{ij}(\chi_i - \chi_j)^2
      \;<\; 1
      \quad \forall\, i.
      \end{equation}
      When $S_i \to 1$ for all $i$ simultaneously, the relaxation rate
      $R_i = c\sqrt{1-S_i} \to 0$: no ordered projected update is possible.
      The network is in the maximally constrained regime of Section~7.1.

    \paragraph{First projectable state.}
      The Big Bang is defined in this framework as the first global
      configuration for which Eq.~\eqref{eq:rc-proj-cond} is satisfied with
      a non-zero spectral gap:
    \begin{equation}
      \label{eq:rc-firstproj}
      \lambda_2\!\left(L_G(t_{\mathrm proj})\right) \;>\; 0,
      \end{equation}
      where $\lambda_2$ is the second eigenvalue of the weighted Laplacian
      $L_G$ of the relaxation network (Appendix~C.1).
      The projectability threshold $\Theta_p$ defined in Appendix~D.5 is
      precisely the condition $\lambda_2 > \varepsilon_p$, so
      $t_{\mathrm proj}$ is the first instant at which $\Theta_p$ is crossed.
      Before $t_{\mathrm proj}$, spacetime, distance, and causal ordering are
      undefined; the apparent singularity of standard cosmology reflects
      only the breakdown of the geometric projection, not a physical
      divergence of $\chi$.

    \paragraph{Energy at first projection.}
      At $t_{\mathrm proj}$, the minimal action available per projectable
      configuration is $\hbar_\chi$, and the total number of simultaneously
      projectable modes is of order unity (by definition of the first
      crossing of $\Theta_p$).
      The effective energy density associated with the first projection is
      therefore set by the Planck energy per Planck volume,
    \begin{equation}
      \label{eq:rc-rho-planck}
      \rho_{\mathrm proj}
      \;\sim\; \frac{E_P}{\ell_P^3}
      \;=\; \frac{c^5}{\hbar G^2}
      \;\approx\; 5 \times 10^{96}~\mathrm{kg\,m^{-3}},
      \end{equation}
      consistently with the standard Planck density.
      The subsequent cascade described in Section~\ref{sec:relax-cascade-hier}
      then accounts for the reduction from $\rho_{\mathrm proj}$ to the
      present-day densities.

    \paragraph{Homogeneity as a structural consequence.}
      Prior to $t_{\mathrm proj}$, all configurations of $\chi$ belong to the
      same fibre of the (not yet defined) projection: they are
      operationally indistinguishable.
      Homogeneity and isotropy of the early Universe are not initial
      conditions to be explained; they are structural consequences of the
      indistinguishability of pre-projection configurations.
      In the language of Appendix~C.2, the horizon problem does not arise
      because causal disconnection is a notion defined only after
      $t_{\mathrm proj}$; the pre-projection substrate is globally connected
      by construction.

    \paragraph{Summary: four linked results.}
      Table~\ref{tab:rc-summary} collects the four results of this section
      and their anchoring equations.

    \begin{table}[h]
      \centering
      \caption{Summary of the four structural results derived in
      Section~\ref{sec:relax-cascade}. Each result is a consequence of
      existing Cosmochrony equations; no new free parameter is introduced.}
      \label{tab:rc-summary}
      \begin{tabular}{lll}
        \hline
        Result & Key relation & Anchoring equation \\
        \hline
        $\Lambda_{\mathrm{eff}}$ dynamical, $w \neq -1$
        & $\rho_\Lambda \propto H^2$
        & Eqs.~(67), (168), \eqref{eq:rc-rho-lambda} \\
        Energy hierarchy $E_P/E_n \sim n$
        & $E_n \sim E_P/n$
        & Eqs.~(202), \eqref{eq:rc-cascade} \\
        $\hbar$ as epoch-dependent invariant
        & $\hbar_\chi = c^3/(K_0\chi_c)$
        & Eqs.~(202), (204), \eqref{eq:rc-hbar-cosmo} \\
        Big Bang $=$ first projectable state
        & $\lambda_2(L_G(t_{\mathrm proj})) > 0$
        & Eqs.~(179), \eqref{eq:rc-firstproj} \\
        \hline
      \end{tabular}
      \end{table}
